% Part: incompleteness
% Chapter: introduction
% Section: overview

\documentclass[../../../include/open-logic-section]{subfiles}

\begin{document}

\olfileid{inc}{int}{ovr}

\olsection{نامکملیت دے نتیجیاں دا اجمالی جائزہ}

ہلبرٹ نوں امید سی کہ ریاضی نوں اک ایہو جیہے !!{axiomatizable}
نظریے وچ رسمی شکل دتی جا سکدی اے جیہدا
!!{complete} تے !!{decidable} ہونا ثابت
کیتا جا سکے۔ اوہ ایہ وی چاہندا سی
کہ بہت کمزور، «متناہی» طریقیاں نال ایس نظریے
دی بے تضادی ثابت ہووے؛ ایس نال وجدانی مکتبِ
فکر دے اعتراضاں دے مقابلے وچ کلاسیکی ریاضی
دا دفاع ہووے گا۔ گوڈل دے نامکملیت دے قضیاں
نے دکھایا کہ ایہ سارے مقصد اکٹھے پورے نہیں
ہو سکدے۔

گوڈل دے پہلے نامکملیت قضیے نے دکھایا کہ
رسل تے وائٹ ہیڈ دی \emph{Principia Mathematica}
دا اک روپ !!{complete} نہیں۔ پر ثبوت خاص
کتاب تک محدود نہیں سی۔ اوہ بہت طرح دے
نظریاں اُتے لاگو ہون والی عام دلیل سی۔
ایس لئی مسئلہ صرف ایہ نہیں کہ
\emph{Principia Mathematica} ریاضی نوں پوری
طرح نہیں سمیٹ سکی، بلکہ ہلبرٹ دے مطلوبہ
خواص رکھن والا \emph{کوئی وی} نظریہ ایہ
کم نہیں کر سکدا۔ ایہ طے کرن وچ کجھ ویلا
لگیا کہ نظریے دیاں کیہڑیاں خاصیتاں نامکملیت
دے قضیاں لئی کافی نیں، تے گوڈل دی دلیل
نوں اینج عام بنایا جائے کہ اوہ صرف ایہناں
خاصیتاں اُتے منحصر ہووے۔ ہن اسی حساب دی
زبان~$\Lang L_A$ وچ نظریاں لئی پہلے
نامکملیت قضیے دا اک بہت عام بیان دے
سکدے آں۔

\begin{thm}
جے $\Gamma$ زبان~$\Lang L_A$ دا اک بے تضاد
تے !!{axiomatizable} نظریہ اے جیہڑا سارے
قابلِ حساب فنکشناں تے سارے !!{decidable}
تعلقّاں دی !!{represents} کردا اے، تاں
$\Gamma$ !!{complete} نہیں۔
\end{thm}

% PNB-CMP-048: move the closing parenthesis after Gamma outside inline math.
% PNB-CMP-049: distinguish ordinary Goedel self-unprovability from Rosser's consistency-only independent witness.
$\Gamma$ دے !!{complete} نہ ہون دا مطلب
ایہ اے کہ گھٹ توں گھٹ اک !!{sentence}~$!A$
ایہو جیہا اے جیہدے لئی $\Gamma \Proves/ !A$ تے
$\Gamma \Proves/ \lnot !A$، دونیں گلاں سچ نیں۔
ایہو جیہے !!a{sentence} نوں نظریے
(یعنی~$\Gamma$) دے لحاظ نال \emph{آزاد}
کہندے نیں۔ آزاد جملیاں دا ہونا اسی نسبتاً
چھیتی ثابت کر سکدے آں۔ گوڈل دی دلیل دی
خاص طاقت ایہ اے کہ اوہ اک \emph{متعین مثال}
بنا کے دیندی اے؛ پر اتھے دو دعویاں نوں
وکھ رکھنا ضروری اے۔ اک !!a{sentence}~$!G_\Gamma$
بنایا جاندا اے، جیہنوں~$\Gamma$ لئی
\emph{گوڈل جملہ} کہندے نیں۔ ایہ
ثابت نہیں ہو سکدا، کیوں جے
$\Gamma$ وچ~$!G_\Gamma$ ایس دعوے دے
برابر اے کہ~$!G_\Gamma$، نظریے
$\Gamma$ وچ ثابت نہیں ہو سکدا۔
صرف بے تضادی توں ایہ نکلدا اے کہ
ایہ جملہ ثابت نہیں ہو سکدا۔
ایس عام گوڈل جملے دی نفی وی ثابت نہ ہون
لئی ہور مضبوط شرط، مثلاً اومیگا-بے تضادی،
چاہیدی اے۔ صرف بے تضادی دی شرط ہیٹھ
اک متعین آزاد !!{sentence} روسر دی بدلی ہوئی
بناوٹ توں ملدا اے۔ گوڈل جملے دی بناوٹ
\emph{تعمیری} اے۔ جے~$\Gamma$ دے مسلّماں
دا بیان تے !!{derivation} دے نظام دی تفصیل
دتی ہووے، ثبوت~$!G_\Gamma$ نوں حقیقت وچ
لکھن دا طریقہ دیندا اے۔

گوڈل دی بناوٹ لئی ضروری اے کہ
$\Lang L_A$ دے اپنے ٹرم تے
!!{formula}s جیہڑیاں خاصیتاں رکھدے نیں، تے
اوہناں اُتے جیہڑے عمل ہوندے نیں، اوہ~$\Lang L_A$ وچ بیان ہو سکن۔
مثلاً «$!A$ اک !!a{sentence} اے»،
«$\delta$،~$!A$ دا اک !!a{derivation} اے»،
تے $\Subst{!A}{t}{x}$ ورگا عمل۔
ایہ طریقہ، اول، نشاناں تے اوہناں دیاں
لڑیواں نوں قدرتی عدداں دے کوڈاں وچ
بدل کے ایہ خاصیتاں تے تعلق بیان کرے۔
ٹرم، !!{formula}s تے !!{derivation}s
ایہو جیہیاں نشان-لڑیاں نیں، جد کہ
$\Lang L_A$ قدرتی عدداں بارے گل
کر سکدی اے۔ دوجا، کوڈ ایویں بنے کہ
$\Gamma$ متعلقہ حقائق ثابت کر سکے۔
ایس لئی دکھانا پینا اے کہ کوڈ شدہ
خاصیتاں قدرتی عدداں دیاں
!!{decidable} خاصیتاں نیں تے
کوڈاں اُتے عمل، قدرتی عدداں اُتے
قابلِ حساب فنکشناں نال ملدے نیں۔
ایہنوں «نحو نوں عددی روپ دینا»
(arithmetization of syntax) کہندے نیں۔

پر نحو نوں عددی روپ دین دی پڑتال
توں پہلے اسی ایہ شرط ویکھاں گے کہ
نظریہ~$\Gamma$ «کافی طاقتور» ہووے،
یعنی سارے قابلِ حساب فنکشناں تے
!!{decidable} تعلقّاں دی
!!{represents} کردا ہووے۔ ایس لئی
«قابلِ حساب» دی ٹھیک تعریف دینی
پینی اے۔ ایہدے کئی طریقے نیں، مثلاً
ٹیورنگ مشین دے نمونے نال، یا کسے
عام مقصد والی پروگرامنگ زبان دے
پروگراماں راہیں حساب کیتے جان والے
فنکشناں نال۔ پر ساڈا مقصد ایہناں
فنکشناں تے تعلقّاں دی زبان~$\Lang L_A$
دے نظریے وچ !!{represents} کرنا اے،
ایس لئی صرف عددی فنکشناں دی
قابلِ حسابیت دی سادی تعریف چُننا
چنگا اے۔ ایہ \emph{بازگشتی فنکشن}
دا تصور اے۔ سو پہلے اسی بازگشتی
فنکشناں نوں ویکھاں گے۔ فیر دکھاواں
گے کہ $\Th{Q}$ پہلے ہی سارے
بازگشتی فنکشناں تے تعلقّاں دی
!!{represents} کردا اے۔ ایس توں
نامکملیت دا قضیہ $\Th{Q}$ تے
$\Th{PA}$ ورگے مخصوص نظریاں
اُتے لاگو کر سکاں گے، کیوں جے
ایہ «کافی طاقتور» نظریاں دیاں
مثالاں ثابت ہو چکیاں ہون گیاں۔
% PNB-CMP-050: remove the frozen source's malformed glossary plural suffix.

نحو نوں عددی روپ دین دا آخری
نتیجہ اک !!a{formula}~$\OProv[\Gamma](x)$
اے۔ ایہ !!{formula}s نوں عدداں
نال کوڈ کر کے~$\Gamma$ دے مسلّماں
توں ثبوت ہون دی گل بیان کردا اے۔
خاص طور تے، جے~$!A$ دا کوڈ
عدد~$n$ اے تے $\Gamma \Proves
!A$، تاں $\Gamma \Proves
\Prov[\Gamma](\num{n})$۔
$\Gamma$ لئی ایہ «ثبوت پذیری
دا محمول» اک خاص مفہوم وچ
$\Gamma$ دی بے تضادی نوں وی
$\Lang L_A$ دے اک !!a{sentence}
وچ بیان کرن دی سہولت دیندا اے:
$\Gamma$ دا «بے تضادی دا جملہ»
!!{sentence}~$\lnot \Prov[\Gamma](\num{n})$
ہووے، جتھے~$n$ کسے تضاد،
مثلاً~$\lfalse$، دا کوڈ اے۔
% PNB-CMP-051: the source's unconditional second-theorem sentence is qualified by its following sentence.
دوجا نامکملیت قضیہ کہندا اے
کہ بے تضاد !!{axiomatizable}
نظریے، جدوں ایس قضیے دیاں
ہور مطلوبہ شرطاں وی پوری کرن،
اپنی بے تضادی دے جملے دا ثبوت
نہیں دے سکدے۔ ایہ شرطاں
صرف سارے قابلِ حساب فنکشناں
تے !!{decidable} تعلقّاں دی
نمائندگی توں کجھ ودھ سخت
نیں؛ اسی دکھاواں گے کہ
$\Th{PA}$ اوہناں نوں پورا
کردا اے۔

\end{document}
